\chapter{System Design and Architecture}
\label{ch:architecture}

\section{Overall System Architecture}
\label{sec:overall_architecture}
The SignBridge system is architected as a decoupled, multi-tier client-server application optimized for real-time video processing and deep learning inference. The high-level architecture comprises two primary tiers:
\begin{enumerate}
    \item \textbf{Presentation Tier (Frontend)}: A single-page web client built with React 19, TypeScript, and Vite that handles user interaction, video file selection, webcam recording via HTML5 media APIs, asynchronous API communication, and dynamic results visualization.
    \item \textbf{Application and Inference Tier (Backend)}: An asynchronous RESTful microservice built with FastAPI and Uvicorn running on Python 3.11, integrating an OpenCV-powered video decoding pipeline and a singleton multi-model neural inference engine caching trained Keras model weights.
\end{enumerate}

Figure \ref{fig:overall_arch} illustrates the comprehensive architectural topology of SignBridge, highlighting the separation of concerns between client-side ingestion, HTTP communication boundaries, backend preprocessing, and the in-memory deep learning model registry.

\begin{figure}[htbp]
\centering
\begin{tikzpicture}[
    node distance=1.2cm and 1.5cm,
    font=\small\sffamily,
    box/.style={draw=blue!70!black, fill=blue!5, rounded corners=4pt, align=center, inner sep=7pt, line width=1pt},
    comp/.style={draw=teal!70!black, fill=teal!5, rounded corners=4pt, align=center, inner sep=7pt, line width=1pt},
    model/.style={draw=purple!70!black, fill=purple!5, rounded corners=4pt, align=center, inner sep=7pt, line width=1pt},
    data/.style={draw=orange!70!black, fill=orange!8, rounded corners=2pt, align=center, inner sep=6pt, line width=0.8pt},
    arrow/.style={->, >=stealth, line width=1.1pt, draw=gray!70!black}
]

% Client Tier
\node[box, text width=4.0cm] (client) {
    \textbf{Client Presentation Tier}\\
    \textit{React 19 + TypeScript + Vite}\\
    \vspace{0.15cm}
    \footnotesize
    $\bullet$ Video Upload (.mov, .mp4, .avi)\\
    $\bullet$ HTML5 Webcam Recorder\\
    $\bullet$ Real-time Latency \& Top-5 UI\\
    $\bullet$ Tri-Model Agreement Matrix
};

% REST Boundaries
\node[data, right=1.6cm of client, text width=2.6cm] (api) {
    \textbf{REST API / HTTP}\\
    \footnotesize
    \texttt{/api/v1/health}\\
    \texttt{/api/v1/model-info}\\
    \texttt{/api/v1/recognize}\\
    \texttt{/api/v1/compare}
};

% Backend Tier
\node[comp, right=1.6cm of api, text width=4.4cm] (backend) {
    \textbf{Backend Processing Tier}\\
    \textit{FastAPI + Uvicorn (Python 3.11)}\\
    \vspace{0.15cm}
    \footnotesize
    $\bullet$ CORS Middleware \& File Validation\\
    $\bullet$ Temporary File Stream Handler\\
    $\bullet$ Multi-Model Lifespan Engine\\
    $\bullet$ Model Registry \& Aliases
};

% Preprocessing Subsystem
\node[comp, below=1.2cm of backend, text width=4.4cm] (preproc) {
    \textbf{OpenCV Video Pipeline}\\
    \textit{app/preprocessing.py}\\
    \vspace{0.15cm}
    \footnotesize
    $\bullet$ Uniform 33-frame Linspace Sample\\
    $\bullet$ BGR $\to$ RGB Color Transform\\
    $\bullet$ Resolution: $224\times224$ / $112\times112$\\
    $\bullet$ Normalization: $[0, 255]$ vs $[0, 1]$
};

% Inference Engine Singleton
\node[model, below=1.2cm of preproc, text width=4.4cm] (engine) {
    \textbf{MultiModelInferenceEngine}\\
    \textit{app/inference.py (Singleton)}\\
    \vspace{0.15cm}
    \footnotesize
    $\bullet$ In-Memory Model Cache\\
    $\bullet$ 76-Class Mapping Resolver\\
    $\bullet$ Top-K Candidate Extraction\\
    $\bullet$ Latency Benchmarking
};

% Loaded Model Artifacts
\node[model, left=1.6cm of engine, text width=4.0cm] (models) {
    \textbf{Trained Keras Model Store}\\
    \textit{backend/artifacts/}\\
    \vspace{0.15cm}
    \footnotesize
    $\bullet$ \textbf{MobileNetV2 Frozen} (67.78\%)\\
    $\bullet$ \textbf{Custom 2D CNN} (5.00\%)\\
    $\bullet$ \textbf{CNN-LSTM Sequence} (4.44\%)\\
    $\bullet$ \textbf{Late Fusion} (60.56\%)
};

% Connectors
\draw[arrow] (client.east) -- node[above, font=\footnotesize] {Upload} (api.west);
\draw[arrow] (api.east) -- (backend.west);
\draw[arrow] (backend.south) -- node[right, font=\footnotesize] {Raw Video} (preproc.north);
\draw[arrow] (preproc.south) -- node[right, font=\footnotesize] {Tensors} (engine.north);
\draw[arrow] (models.east) -- node[above, font=\footnotesize] {Cached Weights} (engine.west);

% Response loop
\draw[arrow, dashed] (engine.south) -- ++(0,-0.4) -| node[near start, below, font=\footnotesize] {Predictions, Top-5, Latency} (client.south);

\end{tikzpicture}

\caption{SignBridge end-to-end system architecture, illustrating client presentation, FastAPI REST boundary, OpenCV preprocessing pipeline, and the cached multi-model inference engine.}
\label{fig:overall_arch}
\end{figure}

\section{Frontend Architecture}
\label{sec:frontend_architecture}
The frontend architecture is organized around modular React functional components, strictly typed TypeScript models, and centralized API communication services. The directory structure is organized as follows:
\begin{itemize}
    \item \texttt{src/components/layout/}: Contains foundational structural components, including the \texttt{Navbar} (which displays dynamic backend health badges and active model indicators) and the \texttt{Footer}.
    \item \texttt{src/components/dashboard/}: Implements the main informational hub, consisting of the \texttt{HeroSection}, the \texttt{ModelCards} component displaying evaluated architectures and benchmark accuracies, the \texttt{PipelineFlow} component, and the \texttt{VocabularyPreview} browser.
    \item \texttt{src/components/recognition/}: Houses the operational recognition workspace, including \texttt{VideoUploader} (drag-and-drop and file selection), \texttt{WebcamRecorder} (camera stream management, countdown timer, and blob generation), \texttt{VideoPreviewCard} (HTML5 video playback), and \texttt{RecognitionResultCard} (displaying predicted ISL words, confidence meters, latency badges, and top-5 probability rankings).
    \item \texttt{src/components/comparison/}: Contains the multi-model comparison suite, including \texttt{BenchmarkSummarySection}, \texttt{ModelComparisonCard}, and \texttt{AgreementAnalysisSection} for evaluating consensus and majority agreement across models.
    \item \texttt{src/services/recognitionApi.ts}: Provides strongly typed asynchronous HTTP client methods (\texttt{analyzeVideo}, \texttt{compareAllModels}, and \texttt{checkBackendHealth}) that communicate with the FastAPI backend.
    \item \texttt{src/types/index.ts}: Defines formal TypeScript interfaces for model metrics, prediction payloads, top-K candidates, and health responses, ensuring end-to-end type safety.
\end{itemize}

\section{Backend Architecture}
\label{sec:backend_architecture}
The backend architecture is structured around FastAPI's asynchronous application lifecycle, implemented across four core Python modules located in \texttt{backend/app/}:
\begin{enumerate}
    \item \texttt{config.py}: Defines global application constants, environment variable bindings, allowed video file extensions (\texttt{.mov}, \texttt{.mp4}, \texttt{.avi}, \texttt{.webm}), maximum upload size limits (100\,MB), frame sampling parameters ($N=33$), and the central \texttt{MODEL\_REGISTRY}. The registry defines metadata, target resolutions, normalizations, and filesystem paths for all registered models.
    \item \texttt{preprocessing.py}: Contains OpenCV video decoding and frame manipulation routines. The core functions \texttt{decode\_video\_frames} and \texttt{preprocess\_for\_model} enforce uniform temporal sampling via \texttt{np.linspace} and apply model-specific resizing and numerical scaling.
    \item \texttt{inference.py}: Implements the \texttt{MultiModelInferenceEngine} class as a global singleton. It manages the in-memory caching of loaded Keras models, loads the verified 76-class label mapping from JSON, computes normalized probability distributions, extracts top-$K$ candidate rankings, and performs equal-weight late fusion.
    \item \texttt{main.py}: Defines the FastAPI application instance, lifespan context manager for startup model pre-loading, CORS middleware configurations, and REST endpoint route handlers.
\end{enumerate}

\section{Model Inference Workflow}
\label{sec:inference_workflow}
The inference workflow translates an incoming raw video stream into a verified lexical sign prediction through five sequential stages, as illustrated in Figure \ref{fig:pipeline_flow}.

\begin{figure}[htbp]
\centering
\begin{tikzpicture}[
    node distance=0.8cm and 0.5cm,
    font=\small\sffamily,
    stepbox/.style={draw=blue!70!black, fill=blue!6, rounded corners=3pt, align=center, inner sep=6pt, text width=2.4cm, line width=0.9pt},
    header/.style={font=\bfseries\footnotesize, text=blue!80!black},
    arrow/.style={->, >=stealth, line width=1.1pt, draw=gray!70!black}
]

\node[stepbox] (step1) {
    \textbf{Stage 01}\\[0.1cm]
    \textbf{Video Ingestion}\\[0.15cm]
    \scriptsize
    $\bullet$ Video Upload / Cam\\
    $\bullet$ Formats: MOV, MP4, AVI\\
    $\bullet$ Size $\le 100$ MB
};

\node[stepbox, right=0.6cm of step1] (step2) {
    \textbf{Stage 02}\\[0.1cm]
    \textbf{Temporal Sampling}\\[0.15cm]
    \scriptsize
    $\bullet$ OpenCV Decode\\
    $\bullet$ 33 Uniform Frames\\
    $\bullet$ \texttt{np.linspace(0, T-1, 33)}
};

\node[stepbox, right=0.6cm of step2] (step3) {
    \textbf{Stage 03}\\[0.1cm]
    \textbf{Spatial Preproc}\\[0.15cm]
    \scriptsize
    $\bullet$ BGR $\to$ RGB\\
    $\bullet$ Resize ($224^2$ / $112^2$)\\
    $\bullet$ Scale: $[0, 255]$ / $[0, 1]$
};

\node[stepbox, right=0.6cm of step3] (step4) {
    \textbf{Stage 04}\\[0.1cm]
    \textbf{Model Inference}\\[0.15cm]
    \scriptsize
    $\bullet$ MobileNetV2 (Frame Avg)\\
    $\bullet$ Custom CNN (Frame Avg)\\
    $\bullet$ CNN-LSTM (Sequence)
};

\node[stepbox, right=0.6cm of step4] (step5) {
    \textbf{Stage 05}\\[0.1cm]
    \textbf{Consensus \& Output}\\[0.15cm]
    \scriptsize
    $\bullet$ Top-1 Predicted ISL Word\\
    $\bullet$ Top-5 Probabilities\\
    $\bullet$ Cross-Model Agreement
};

\draw[arrow] (step1) -- (step2);
\draw[arrow] (step2) -- (step3);
\draw[arrow] (step3) -- (step4);
\draw[arrow] (step4) -- (step5);

\end{tikzpicture}

\caption{Five-stage inference pipeline flow in SignBridge, detailing ingestion, temporal linspace sampling, spatial standardization, model inference, and consensus output.}
\label{fig:pipeline_flow}
\end{figure}

The detailed operational mechanics of each stage are:
\begin{itemize}
    \item \textbf{Stage 1 (Video Ingestion)}: The user uploads an RGB video file or captures a clip via webcam. The backend validates the file extension and streams chunks into a named temporary file on disk while enforcing a 100\,MB ceiling.
    \item \textbf{Stage 2 (Uniform Frame Sampling)}: OpenCV opens the video container via \texttt{cv2.VideoCapture}. If the total frame count $T$ is readable, indices are generated via:
    \begin{equation}
        \mathcal{I} = \left\lfloor \text{linspace}(0, T - 1, 33) \right\rfloor \in \mathbb{Z}^{33}
        \label{eq:linspace}
    \end{equation}
    The capture pointer seeks directly to each index $\tau \in \mathcal{I}$, decoding the corresponding frame. If container metadata is missing, sequential decoding with linspace subsampling is executed as a fallback.
    \item \textbf{Stage 3 (Spatial Preprocessing)}: Each decoded frame is converted from BGR to RGB color-space and resized using \texttt{cv2.INTER\_AREA} interpolation to $224\times224$ (for MobileNetV2 and Custom CNN) or $112\times112$ (for CNN-LSTM). Numerical values are cast to float32 and scaled appropriately.
    \item \textbf{Stage 4 (Multi-Model Inference)}: The preprocessed tensor is routed to the selected model in the cached inference engine. Frame-level models compute forward passes across all 33 frames, while the recurrent CNN-LSTM evaluates the 5D sequence tensor directly.
    \item \textbf{Stage 5 (Consensus \& Aggregation)}: For frame-level architectures, the 33 frame probability vectors are averaged. The top predicted class index $\hat{k} = \arg\max_k \bar{p}_k$ is mapped to the corresponding ISL word label, top-5 candidates are extracted, and cross-model agreement metrics are computed.
\end{itemize}

\section{API Communication and Protocols}
\label{sec:api_communication}
Communication between the React client and FastAPI backend occurs over standard HTTP/1.1 protocols using JSON and Multipart Form-Data encodings. Table \ref{tab:api_endpoints_summary} outlines the primary API endpoints exposed by the SignBridge backend.

\begin{table}[htbp]
\centering
\caption{SignBridge REST API endpoint specifications.}
\label{tab:api_endpoints_summary}
\small
\begin{tabularx}{\textwidth}{l l p{4.2cm} X}
\toprule
\textbf{HTTP Method} & \textbf{Endpoint Path} & \textbf{Request Payload} & \textbf{Response Description} \\
\midrule
\texttt{GET} & \texttt{/api/v1/health} & None & Returns service health status, class count (76), default model ID, and load status per registered model. \\
\addlinespace
\texttt{GET} & \texttt{/api/v1/model-info} & Optional query: \texttt{?model=<id>} & Returns architecture specifications, input/output tensor shapes, and preprocessing parameters. \\
\addlinespace
\texttt{POST} & \texttt{/api/v1/recognize} & \texttt{multipart/form-data}: \texttt{video} (File), \texttt{model} (Form/Query string) & Decodes video, runs inference on specified model, and returns top-1 word, confidence, top-5 candidates, and latency. \\
\addlinespace
\texttt{POST} & \texttt{/api/v1/compare} & \texttt{multipart/form-data}: \texttt{video} (File) & Evaluates all three models in parallel on the same clip; returns individual predictions, latencies, and consensus metrics. \\
\bottomrule
\end{tabularx}
\end{table}

\section{Data Flow}
\label{sec:data_flow}
Data traverses the SignBridge pipeline through tightly defined representations:
\begin{enumerate}
    \item \textbf{Binary Video Bitstream}: Raw compressed video packets (\texttt{.MOV}, \texttt{.mp4}, \texttt{.webm}) arrive from the browser and are written to a temporary storage path on the backend host.
    \item \textbf{Decoded Raw Frame List}: A Python list of 33 NumPy arrays of shape $(H, W, 3)$ with \texttt{uint8} values in $[0, 255]$ in RGB color order.
    \item \textbf{Preprocessed Model Tensor}:
    \begin{itemize}
        \item For MobileNetV2: 4D array of shape $(33, 224, 224, 3)$ of type \texttt{float32} in range $[0, 255]$.
        \item For Custom CNN: 4D array of shape $(33, 224, 224, 3)$ of type \texttt{float32} in range $[0, 1]$.
        \item For CNN-LSTM: 5D array of shape $(1, 33, 112, 112, 3)$ of type \texttt{float32} in range $[0, 1]$.
    \end{itemize}
    \item \textbf{Logit and Softmax Distributions}: Raw neural output tensors transformed by the terminal Softmax activation into a discrete probability distribution over 76 classes, satisfying $\sum_{k=1}^{76} p_k = 1.0$.
    \item \textbf{Structured JSON Payload}: Serialized JSON response containing predicted lexical words, floating-point confidence values, top-$K$ candidate arrays, and execution timing metrics.
\end{enumerate}

\section{System Workflow}
\label{sec:system_workflow}
Figure \ref{fig:workflow_frontend} illustrates the complete user interaction workflow implemented in the frontend application, showing validation checkpoints, mode selection, and error-handling pathways.

\begin{figure}[htbp]
\centering
\begin{tikzpicture}[
    node distance=0.8cm and 1.2cm,
    font=\small\sffamily,
    action/.style={draw=blue!70!black, fill=blue!5, rounded corners=3pt, align=center, inner sep=5pt, text width=3.8cm, line width=0.8pt},
    decision/.style={draw=orange!80!black, fill=orange!8, diamond, aspect=1.8, align=center, inner sep=2pt, font=\footnotesize\sffamily, line width=0.8pt},
    service/.style={draw=teal!70!black, fill=teal!8, rounded corners=3pt, align=center, inner sep=5pt, text width=3.8cm, line width=0.8pt},
    ui/.style={draw=purple!70!black, fill=purple!6, rounded corners=3pt, align=center, inner sep=5pt, text width=3.8cm, line width=0.8pt},
    arrow/.style={->, >=stealth, line width=1pt, draw=gray!70!black}
]

\node[action] (input) {
    \textbf{User Interaction}\\
    Upload Video Clip (.mov, .mp4)\\
    \textit{or} Record Live Webcam Clip
};

\node[decision, below=0.6cm of input] (val) {
    File Valid?\\
    $\le 100$ MB?
};

\node[action, right=1.2cm of val, text width=2.6cm, draw=red!70!black, fill=red!6] (err) {
    \footnotesize Display Toast / Banner Error
};

\node[action, below=0.7cm of val] (mode) {
    \textbf{Select Operation}\\
    $\bullet$ Single Model Recognition\\
    $\bullet$ Multi-Model Comparison
};

\node[service, below=0.7cm of mode] (api) {
    \textbf{FastAPI REST Request}\\
    \texttt{POST /api/v1/recognize}\\
    \textit{or} \texttt{POST /api/v1/compare}
};

\node[ui, below=0.7cm of api] (result) {
    \textbf{Interactive Results View}\\
    $\bullet$ Primary Recognized ISL Word\\
    $\bullet$ Top-5 Confidence Distribution\\
    $\bullet$ Measured Latency (ms)\\
    $\bullet$ Cross-Model Consensus
};

\draw[arrow] (input) -- (val);
\draw[arrow] (val) -- node[above, font=\footnotesize] {No} (err);
\draw[arrow] (val) -- node[right, font=\footnotesize] {Yes} (mode);
\draw[arrow] (mode) -- (api);
\draw[arrow] (api) -- (result);

\end{tikzpicture}

\caption{SignBridge frontend user interaction and execution workflow, from file selection or webcam capture through client validation, REST dispatch, and results visualization.}
\label{fig:workflow_frontend}
\end{figure}

The user may either drag-and-drop an existing sign video clip or record a live clip using their webcam. The client performs instant format and size verification; if valid, the user chooses between single-model recognition or the tri-model comparison suite. Upon receiving the JSON response, the frontend renders interactive confidence meters, candidate lists, and agreement banners.
